\documentclass[10.5pt,a4paper]{article}
\usepackage[margin=2cm]{geometry}
\usepackage{amsmath,amssymb}
\usepackage{booktabs,array,tabularx}
\usepackage{enumitem}
\usepackage[dvipsnames]{xcolor}
\usepackage{hyperref}
\usepackage{titlesec}
\usepackage{mdframed}
\hypersetup{colorlinks=true,linkcolor=NavyBlue,urlcolor=NavyBlue}
\setlist{itemsep=1pt,topsep=3pt}
\setlength{\parindent}{0pt}
\setlength{\parskip}{4pt}
\titleformat{\section}{\Large\bfseries\color{NavyBlue}}{\thesection}{0.6em}{}
\titleformat{\subsection}{\large\bfseries}{\thesubsection}{0.6em}{}
\newcommand{\Cu}[1]{#1$\times$#1$\times$#1}
\newcommand{\Sq}[1]{#1$\times$#1$\times$1}
\newcommand{\M}[1]{\textsc{#1}}
\newmdenv[linecolor=NavyBlue,linewidth=0.8pt,backgroundcolor=NavyBlue!4,innertopmargin=4pt,innerbottommargin=4pt]{keybox}
\newmdenv[linecolor=BrickRed,linewidth=0.8pt,backgroundcolor=BrickRed!4,innertopmargin=4pt,innerbottommargin=4pt]{warnbox}
\newcommand{\sub}[1]{\par\smallskip\textbf{#1}\quad}

\title{\textbf{Study guide: how every method works}\\[2pt]\large Short Solutions for Large Rubik's Cubes and Squares}
\author{Companion to \texttt{notes.pdf} (full detail and question bank)}
\date{18 November 2026}

\begin{document}
\maketitle
\tableofcontents

\section{The story in 30 seconds}
\begin{keybox}
``Finding the shortest solution of a big Rubik's cube is NP-complete, so in practice people use fast solvers that find
\emph{short} solutions. We built one engine for cubes and Squares of any size, and ran 8 solvers plus 3 outside baselines on
the same puzzles. Our main method, \M{SymPort}, beats the best published Kaggle Santa 2023 solutions on average for cubes
$N=4$--7. \M{NetGain} is the best single solver on large random cubes. A counting lower bound shows everyone is still about
$2\times$ above optimal, and that gap grows like $\log N$.''
\end{keybox}
Four research questions: covering every size; which method is shortest at each size; can anyone beat the published
solutions; how far is everyone from optimal.

\textbf{God's number} = the worst case: the fewest moves that solve \emph{every} state (the diameter of the state graph).
\Cu{3}: 20 face turns (Rokicki et al.\ 2010), 26 quarter turns (2014). \Cu{N}: grows as $\Theta(N^2/\log N)$ (Demaine et
al.\ 2011). Our exact Square values: 6, 10, 12 for $N=3,4,5$.

\textbf{NP-complete} (Demaine, Eisenstat, Rudoy 2018): ``can this state be solved in at most $k$ moves?'' is NP-complete
for the \Cu{N}. Any solution is easy to find (reduction does it in polynomial time); only the \emph{shortest} is hard.
So beyond small sizes we build fast solvers and measure them against bounds.

\section{Building blocks}
\subsection{The two puzzles}
\begin{itemize}
\item \textbf{Cube} \Cu{N}: $6N^2$ stickers, $3N$ slices, each turned $90^\circ$, $180^\circ$ or $270^\circ$: $9N$ moves.
\item \textbf{Square} \Sq{N}: a flat grid; a move flips a whole row or column $180^\circ$: $2N$ moves, each its own inverse.
\item \textbf{Why the Square:} same internal structure as the cube, but exactly solvable up to $N=5$. It is our laboratory for
optimality and for training networks cheaply.
\end{itemize}

\subsection{Clusters (behind almost every method)}
A piece only ever moves among positions of its own kind; each closed group is a \textbf{cluster}. Solving the puzzle = solving
every cluster. Cube: 1 corner cluster, $\lfloor N/2\rfloor-1$ wing clusters, $\approx(N-2)^2/4$ center clusters of 24 stickers
(the \Cu{20}: 9 wing, 81 center). The count grows like $N^2$, but each cluster is a small fixed-size problem.

\textbf{Implementation} (\texttt{engine/orbits.py}): union-find over the move permutations gives the orbits; stickers in a
cluster are sorted in an order that does \emph{not} depend on $N$. A center cluster on a \Cu{5} looks identical in code to one
on a \Cu{9}. That is what lets one network (\M{OneNet}) work at every size.

\subsection{Commutators and 3-cycles}
A commutator is $[A,B]=A\,B\,A^{-1}B^{-1}$. If $A$ and $B$ barely overlap, most of $A$ is undone by $A^{-1}$; only pieces at the
overlap move. Chosen well, exactly \textbf{3 pieces of one cluster} rotate and nothing else moves: a \emph{pure 3-cycle}, the
smallest even permutation, enough to solve any cluster.
\begin{itemize}
\item Cube: forms like $[A,\,B\,C\,B^{-1}]$ (inner slice, outer face, inner slice).
\item Any target cycle: conjugate a base one by a setup, $s\cdot\text{cycle}\cdot s^{-1}$ (still pure). BFS over setups of up to
4 moves reaches all $24\cdot23\cdot22/3=4{,}048$ 3-cycles of a 24-sticker cluster.
\item Square: $c\,r\,c\,r$ changes only the cell at the crossing and its two mirror cells (Demaine's Lemma 2). Checked in the
engine: on a \Sq{5}, column 3 with row 1 moves cells (1,3), (1,1), (3,3).
\end{itemize}

\subsection{Metrics and goals}
\textbf{HTM}: any turn of a slice = 1 (random states). \textbf{QTM}: a half turn = 2 (\textbf{the Santa score is exactly
QTM}). \textbf{Canonical form}: same-axis moves commute, so they are merged and sorted; it prunes searches and defines the
counting bound. Goals: random benchmark = solved in any orientation (24 for the cube); Santa = one \emph{exact} state.

\subsection{The harness: why every number is trustworthy}
Every solver returns a move list; \textbf{our engine re-applies it and decides ``solved'' itself}. Crash or timeout = failure,
kept in the denominator. Each result row stores seed, hardware and git commit; every table and figure is regenerated by
\texttt{paper/scripts/make\_tables.py}.

\section{The methods}
Three families: \textbf{exact search} (\M{IDA*-PDB}); \textbf{cluster schedules} (\M{Reduction}, \M{Demaine}, \M{NetGain});
\textbf{learned and portfolio} (\M{CayleyPy}, \M{OneNet}, \M{KMC-Repro}, \M{SymPort}).

\subsection{\M{IDA*-PDB} (m1): optimal, small puzzles only}
\sub{Idea.} Depth-first search with an increasing limit; prune when moves so far + lower-bound estimate of moves left exceed it.
If the estimate never overestimates (admissible), the first solution found is optimal.
\sub{Estimate.} Pattern databases: exact distances of sub-problems, e.g.\ the 8 corners (88 million entries) and two sets of 6
edges; $h=\max$ of the lookups. A sub-problem never needs more moves than the whole puzzle, so $h$ is admissible. Square: max
over clusters of each cluster's exact distance.
\sub{Implementation.} Tables by BFS at 4 bits/entry; Numba search, parallel over the first move, $\approx1.2\cdot10^7$ nodes/s;
canonical pruning. BFS from solved gave the exact Square God's numbers 6, 10, 12.
\sub{Result.} Optimal on every random Square $N\le5$ (\Sq{5} $\approx$13\,s); none at $N=6$. \Cu{3}: 91/100 twelve-move scrambles.
\sub{Why it stops.} Cost $\approx b^{\,d-h}$, $b\approx13$: exponential. Korf's 1997 runs used days per instance; we capped at
30 minutes. Our PDB sizes match his: a budget limit, not a bug.

\subsection{\M{Reduction} (m2): how humans solve big cubes}
\sub{Order.} Wing parity (one slice turn per odd wing cluster, since 3-cycles are even) $\to$ centers by greedy 3-cycles (most
pieces fixed per move) $\to$ corners and middle edges as a \Cu{3} by Kociemba's two-phase $\to$ wings by 3-cycles. Outer turns
keep centers on their faces and pure 3-cycles touch nothing else, so stages never break each other.
\sub{Square.} Frame solved exactly on its projection; each interior cluster with $c\,r\,c\,r$, at most 3 per cluster.
\sub{Weakness.} Each 3-cycle fixes $O(1)$ pieces: $\Theta(N^2)$ moves (1,885 HTM on a random \Cu{10}).

\subsection{\M{Demaine} (m3): the optimal growth rate}
\sub{Idea.} Do the \emph{same} fix to many clusters with one shared sequence: turning column set $X$ and row set $Y$ as
$C_X R_Y C_X R_Y$ applies the commutator to every cluster in $X\times Y$ for only $2|X|+2|Y|$ moves. This gives
$\Theta(N^2/\log N)$.
\sub{Square.} Faithful: each cluster's optimal commutator word from its table; round $t$ applies everyone's $t$-th commutator in
batched blocks.
\sub{Cube (adapted).} Demaine batches clusters in identical states, which essentially never occur at practical sizes. We batch
\textbf{3-cycle types} instead: one template $[A_i,\,B\,C_j\,B^{-1}]$ applied to a block $I\times J$ for $2|I|+2|J|+4$ moves,
used only if it \textbf{helps every cluster in it and harms none}.
\sub{Weakness.} Unanimous blocks shrink to single cells: on the \Cu{20}, 244 of 292 center rounds fixed only one cluster.

\subsection{\M{NetGain} (m5c, new): the ``block as a whole pays'' rule}
\sub{Idea.} Drop unanimity; apply a block when its total gain beats the best single fix per move:
\[(X^\ast,Y^\ast)=\arg\max\sum_{(i,j)\in X\times Y} g(i,j)\;-\;\lambda\,(\text{block cost}),\]
$\lambda$ = best pieces/move of any single-cluster 3-cycle now. Some cells may gain 0 or even lose.
\sub{Implementation.} Alternating maximisation (rows given columns, columns given rows) from the 3 best rows; each candidate
scored exactly by applying it; setup sharing (a repeated setup's undo and redo cancel, charged once); one-step lookahead for
$N\le15$. Everything else is copied from \M{Demaine}, so the difference is the schedule alone: a clean experiment.
\sub{Result.} \Cu{20}: 229 of 333 rounds bulk; \textbf{4,925 vs 6,612 HTM}. 20 random cubes per size, paired 95\% CIs: shorter
from $N=5$, \textbf{$-17.2\%$ at $N=10$}. Best single solver on large Squares.
\sub{Limits.} No proved growth bound; weak on Santa (exact goal, QTM): 2.8--3.3$\times$ published.

\subsection{\M{CayleyPy} (m4): learned distance + beam search}
\sub{Training.} Random walks from solved that never undo the last move; each state labelled with its step $k$ as a proxy for
distance; fresh data every epoch. ResMLP on one-hot stickers, MSE on $k$, Adam $10^{-3}$, batch 1024. Square models 0.25--0.45M
parameters, 1--4 CPU-minutes each; a larger GPU model (Colab T4) at $N=8$.
\sub{Search.} Beam: keep the $W$ best states by network score, expand every move, dedup by 64-bit hash, repeat.
\sub{Result.} Optimal on all \Sq{3}, \Sq{4}; at the optimum at $N=5$; best (with \M{SymPort}) up to $N=8$. Wider beams help
(\Sq{8}: 3/10 at $W=16$, all from $W=1024$).
\sub{No cubes.} One network per size, billions of examples ($\approx$8B for their \Cu{3}): beyond our compute.

\subsection{\M{OneNet} (m5a, new): one network for every size}
\sub{Idea.} One \textbf{token per cluster}. Cluster order is $N$-invariant, so a token means the same at every size; a bigger
cube is just more tokens.
\sub{Model.} Token = cluster colours + cluster-kind embedding + position features (depth$/(N-1)$, $N/20$). Transformer with
\textbf{no positional encoding} (clusters are a set). \textbf{Additive head}: each token predicts its cost and they are summed:
extends to any token count and gives per-cluster costs for staged search.
\sub{Training.} Walks on $N=4$--8; loss normalised by $N^2/\ln N$; AdamW, OneCycle. Main model $d=192$, 4 layers, 1.56M
parameters on a T4.
\sub{Search.} $N\le7$: full beam. $N\ge8$: staged outside-in (frame, edges, interiors), each stage an exact goal so later stages
cannot undo earlier ones; $c\,r\,c\,r$ macros; small stages solved exactly by BFS.
\sub{Result.} Solves unseen sizes (up to \Sq{14}); $r=0.66/0.65/0.53$ at unseen $N=10/14/20$; second only to per-size networks
up to $N=7$. Honest: at large $N$ the stage structure carries much of the load.

\subsection{\M{KMC-Repro} (m6): the winners, rebuilt from their write-up}
\sub{Insertion identity.} Store the permutation of every suffix of the solution; inserting $B$ before move $t$ changes the final
state by $C=\mathrm{inv}_t\circ B\circ\mathrm{suf}_t$: one lookup, no replay, so millions of edits are cheap to score.
\sub{Pipeline.} Special parts (corners as an optimal \Cu{2}, or the central \Cu{3}) $\to$ wing parity $\to$ simulated annealing
(QTM $+\,2.5\times$ wrong stickers) $\to$ greedy 4-move commutator insertion $\to$ optimal 3-cycles.
\sub{Result.} 1.61--1.71$\times$ published; one run of their real program is 1.05--1.18$\times$. The gap measures what the write-up
leaves out (bulk commutators, exact 3-cycle counting, big beam, optimal central stage). Built from the write-up only, before
reading their code.

\subsection{\M{SymPort} (m7, new): the main method}
\sub{Idea.} The cube has 48 symmetries (24 rotations $\times$ mirror). A rotated/mirrored and recoloured scramble is a different
but equivalent puzzle. Solve $k$ copies with KMCoders' program, map each answer back, \textbf{keep the shortest}. Seed-free
diversity. The Square has 8 symmetries.

\begin{keybox}
\textbf{Proof (board-ready).} A rotation or reflection $R$ maps every slab onto a slab, so $R\,q\,R^{-1}$ is a turn
$\phi_R(q)$ of another slab by the same angle: the \textbf{same cost}. If $q_1\ldots q_L$ solves the copy, then
$\phi_R(q_1)\ldots\phi_R(q_L)$ solves the original: substitute $q_i=R^{-1}\phi_R(q_i)R$ and the inner $R^{-1}R$ pairs
telescope. Same length in QTM, HTM and STM. The harness re-verifies every mapped-back solution anyway.
\end{keybox}

\sub{Configuration.} KMCoders' beam variant; 1M annealing iterations instead of 20M (no measurable loss); odd $N$: optimal
central \Cu{3} from Rokicki's nxopt. $k=48$ at $N=4$, 24 at 5--7, 12 at 8--9.
\sub{Results (mean QTM).}
\begin{center}\small
\begin{tabular}{@{}lrrrrrr@{}}
\toprule
$N$ & 4 & 5 & 6 & 7 & 8 & 9 \\
\midrule
\M{SymPort} & \textbf{64.1} & \textbf{95.0} & \textbf{158.3} & \textbf{205.2} & 294.7 & 361.3 \\
Published & 66.7 & 97.1 & 161.0 & 208.8 & \textbf{291.3} & \textbf{354.7} \\
\bottomrule
\end{tabular}
\end{center}
$N=4$--7: 6,334 vs 6,514 ($-2.8\%$), W/T/L 41/14/17. Best at all 18 Square sizes.

\sub{Why it loses at $N=8$, 9 (own this).} \textbf{gain} = (median copy $-$ best copy)/median; \textbf{gap} = (median copy $-$
published)/published. \M{SymPort} wins exactly where gain $>$ gap.
\begin{enumerate}
\item \textbf{The spread does not grow:} copy SD stays $\approx$4.5--7.4 QTM while lengths grow like $N^2$; relative spread
falls 6\% $\to$ 2\%. Copies differ early; the long late stages come out alike.
\item \textbf{Fewer copies:} 31--47 CPU-min each, so only 12.
\item \textbf{Published is itself a portfolio:} best of many runs over the competition.
\end{enumerate}
Projection (normal model, matches the measured best-of-12): 48 copies give 290.2 vs 291.3 at $N=8$ (just wins), 357.3 vs 354.7
at $N=9$ (still loses). Two consecutive losses with a widening gap, plus this projection, are why we stopped at $N=9$.

\begin{warnbox}
\textbf{Be honest about:} the seed-control experiment was \textbf{not run} (cannot claim symmetries beat 48 random seeds);
3--4 puzzles per size at $N\ge6$ (margins not significant); copies killed by the OS for lack of memory were rerun, every
finished copy was kept (not cherry-picking).
\end{warnbox}

\section{The counting lower bound}
\sub{Argument.} Every state at distance $d$ is reached by at least one canonical word of length $d$. If short words are fewer
than states, some state is far: the smallest $L$ with $\sum_{d\le L}\text{words}(d)\ge|G|$ is a lower bound on God's number.
\sub{Counting.} Cube: per layer $(1+x)^2$, block $f(x)=(1+x)^{2N}-1$, axes alternate, growth
$b=\bigl((3/2)^{1/(2N)}-1\bigr)^{-1}$. Square: $b=1/(2^{1/N}-1)$.
\sub{Why $N^2/\log N$.} $\ln|G|\sim N^2$ (number of clusters), $\ln b\sim\log N$.
\sub{Calibration.} Square bound / true God's number: 4/6, 6/10, 8/12 (ratio 1.5--1.7).
\sub{Finding.} Published Santa lengths are 1.81$\times$ the bound at $N=4$, 2.22$\times$ at $N=10$, a gap growing like $\log N$.
Their final stage fixes $O(1)$ pieces per 3-cycle ($N^2$ growth); Demaine-style batching in the endgame (what \M{NetGain} does)
is the way to close it.

\section{Suggested talk flow (12--15 min)}
\begin{enumerate}
\item Problem: puzzles, sizes, NP-complete, God's number (infographic slides 6--8).
\item Engine: clusters, commutators, harness re-verification (the ``trust'' slide).
\item Methods by family: exact $\to$ cluster schedules (\M{Reduction} $\to$ \M{Demaine} $\to$ \M{NetGain} as a progression) $\to$
learned (\M{CayleyPy} $\to$ \M{OneNet}).
\item \M{SymPort}: idea, proof, Table 2.
\item Why it stops at $N=8$: gain vs gap. Examiners like an honest limit.
\item Bound: everyone is $\approx2\times$ optimal and the gap grows.
\item Future work: seed control, recombining the best stages across copies, adaptive budget, stronger base solver.
\item Demo (\texttt{demo/README.md} run-through).
\end{enumerate}

\section{Five traps to pre-empt}
\begin{warnbox}
\begin{itemize}
\item \textbf{21\% vs 17.2\% at $N=10$:} abstract's 21\% is from the 2-state table; the 20-state CI table says 17.2\%. Name the
source of each.
\item \textbf{Slide 13 says OneNet $r\ge0.70$:} the paper says 0.66 / 0.65 / 0.53. Fix the slide or correct it aloud.
\item \textbf{``Isn't \M{SymPort} just running their solver 48 times?''} The inner solver is theirs, and we say so. The
contributions are the symmetry-diversity idea with its exact-length proof, the configuration, and the measured limits.
\item \textbf{Never say ``optimal'':} say ``best up to $N=7$''; everyone is $\approx2\times$ the bound.
\item \textbf{\texttt{m5b\_stages}} in the code: a dropped staged variant, not in the paper.
\end{itemize}
\end{warnbox}

\end{document}
